Ce travail pratique s'inscrit dans le domaine de l'optique géométrique,
plus précisément dans l'étude des systèmes centrés minces.
Il a pour objectif de déterminer les \textbf{distances focales
de deux lentilles}, l'une convergente et l'autre divergente,
en utilisant et en comparant \textbf{différentes méthodes focométriques}
telles que les méthodes de Bessel, Silbermann, Badal
ou encore celle des points conjugués \textbf{en tenant compte des incertitudes.}

La problématique posée est la suivante :
\textbf{dans quelle mesure les différentes techniques de mesure expérimentale
permettent-elles d'obtenir une valeur précise de la distance focale},
et dans une moindre mesure, comment la confrontation des incertitudes absolues et relatives 
permet-elle de \textbf{hiérarchiser la fiabilité de ces méthodes} ?